\documentclass[11pt,reqno]{amsart}
\usepackage[T1]{fontenc}
\usepackage{lmodern}
\usepackage[margin=1in]{geometry}
\usepackage{amsmath,amssymb,mathtools,booktabs,microtype}
\usepackage[hidelinks]{hyperref}
\hypersetup{pdftitle={Galois-averaged indicators and sphericality in rank three},pdfauthor={}}
\newtheorem{theorem}{Theorem}[section]
\newtheorem{proposition}[theorem]{Proposition}
\newtheorem{lemma}[theorem]{Lemma}
\newtheorem{corollary}[theorem]{Corollary}
\theoremstyle{definition}
\newtheorem{definition}[theorem]{Definition}
\theoremstyle{remark}
\newtheorem{remark}[theorem]{Remark}
\newcommand{\C}{\mathbb C}
\newcommand{\Q}{\mathbb Q}
\newcommand{\Z}{\mathbb Z}
\newcommand{\T}{\mathbb T}
\newcommand{\cC}{\mathcal C}
\newcommand{\cS}{\mathcal S}
\newcommand{\cD}{\mathcal D}
\newcommand{\one}{\mathbf 1}
\newcommand{\eps}{\varepsilon}
\newcommand{\Irr}{\operatorname{Irr}}
\newcommand{\Tr}{\operatorname{Tr}}
\newcommand{\tr}{\operatorname{tr}}
\newcommand{\FPdim}{\operatorname{FPdim}}
\newcommand{\Gal}{\operatorname{Gal}}
\newcommand{\Rep}{\operatorname{Rep}}
\newcommand{\Vect}{\operatorname{Vec}}
\newcommand{\Hom}{\operatorname{Hom}}
\newcommand{\id}{\operatorname{id}}
\newcommand{\Dirac}{\operatorname{Dirac}}
\numberwithin{equation}{section}
\title[Galois-averaged indicators in rank three]{Galois-averaged indicators and sphericality in rank three}
\subjclass[2020]{18M20, 18M15, 16T05}
\keywords{Fusion category, spherical structure, formal codegree, Frobenius--Schur indicator, Galois averaging}
\begin{document}
\begin{abstract}
We prove that every complex fusion category of rank three admits a spherical
structure. The cubic case follows from Galois averaging of six even
Frobenius--Schur indicators in the sphericalization: a positive measure on
central twists and a short trigonometric inequality leave only the
quantum $\mathfrak{sl}_2$ fusion ring. We prove the quadratic case and
classify the Grothendieck rings directly, without a multiplicity bound.
\end{abstract}
\maketitle
\paragraph{Research status.} \textbf{AI-generated research \textperiodcentered{} independent verification pending.} Treat this manuscript and its accompanying research data with caution: they have not been independently verified by a human mathematician or formalized in Lean.

\section{Introduction}

A pivotal structure on a fusion category is a monoidal natural isomorphism from the identity functor to the double-dual functor. It is spherical when left and right categorical traces agree. Etingof, Nikshych and Ostrik conjectured that every fusion category admits a pivotal structure \cite[Conjecture~2.8]{ENO}. Ostrik classified rank-three fusion categories under this hypothesis \cite[Theorem~1.1]{Ostrik}. In the same paper he established sphericality for the rational and quadratic cases, but left a cubic case without a multiplicity bound \cite[Section~5]{Ostrik}.

\begin{theorem}\label{thm:main}
Every fusion category of rank three over $\C$ admits a spherical structure.
\end{theorem}

Thus Ostrik's tensor-equivalence classification, previously stated under
pivotality, holds without that assumption. We give a separate direct proof
of the five possible fusion rings in Corollary~\ref{cor:list}. The cubic
sphericality proof is the new argument in this paper, and we give the
rational and quadratic sphericality proofs as well.

The main point is to average \emph{normalized} twist moments over the dimension field of the sphericalization. This field is cubic, whereas the indicators of orders $2,4,6,8,10,12$ lie in $\Q(\zeta_{120})$, a field of degree $32$. All three embeddings of the dimension field therefore extend simultaneously while fixing these indicators. A Frobenius-algebra trace identity annihilates every non-unit term in the average. Positivity comes from squares of real central dimensions, not from pseudo-unitarity. This yields a uniform obstruction rather than a search over signed fusion tables.

Throughout, a fusion category is finite, semisimple, rigid and $\C$-linear, with simple tensor unit. Its rank is the number of isomorphism classes of simple objects. We write $K_0(\cC)$ for its Grothendieck ring and $Z(\cC)$ for its Drinfeld center.

\section{Rank-three rings and their sphericalizations}

\subsection{The two codegree coefficients}

A rank-three fusion ring is either $\Z[C_3]$, or all its basis elements are self-dual and it has the form $K(k,l,m,n)$, with basis $1,X,Y$ and multiplication
\begin{equation}\label{eq:ring}
 X^2=1+mX+kY,\qquad Y^2=1+lX+nY,\qquad XY=kX+lY.
\end{equation}
Here $k,l,m,n$ are nonnegative integers satisfying
\begin{equation}\label{eq:assoc}
 k^2+l^2=kn+lm+1.
\end{equation}
Here is a quick derivation. If $X^*=Y\ne X$, the product $X^2$ has no
unit constituent; writing $d=\FPdim(X)=\FPdim(Y)$ gives
$d^2=(a+b)d$ for nonnegative integers $a,b$, so $d$ is an integer.
The product $XY$ has one unit constituent, hence
$d^2=1+(c+e)d$ for integers $c,e\geq0$; therefore $d\mid1$ and
$d=1$, giving $\Z[C_3]$. Otherwise all three simples are self-dual.
The involution on a based ring reverses products, so the ring is
commutative. Frobenius reciprocity forces the coefficients in $XY$
to be $k,l$ as displayed; associativity of $X^2Y=XYX$ gives
\eqref{eq:assoc}, which also suffices for associativity.

Let $f_1,f_2,f_3$ be the formal codegrees. They are the eigenvalues of multiplication by
\[
 \Omega=1+X^2+Y^2=3+(l+m)X+(k+n)Y.
\]
They are positive, and their reciprocal sum is one. Consequently their characteristic polynomial has the form
\begin{equation}\label{eq:poly}
 p(t)=\prod_{i=1}^3(t-f_i)=t^3-bt^2+ct-c.
\end{equation}
We call the ring \emph{cubic} when $p$ is irreducible over $\Q$. We will use the following rank-three identities:
\begin{align}
 b&=9+(l+m)^2+(k+n)^2
   =7+3k^2+3l^2+m^2+n^2,\label{eq:b}\\
 c&=4b-9+(kl-mn)^2+\bigl(2(k^2-lm)-1\bigr)^2.\label{eq:c}
\end{align}
The first is the trace of multiplication by $\Omega$, followed by \eqref{eq:assoc}. The second is its determinant, simplified using the same relation; it is also \cite[equation~(12)]{Ostrik}. In particular,
\begin{equation}\label{eq:lower}
 c\geq 4b-8,
\end{equation}
because the last square in \eqref{eq:c} is the square of an odd integer. Vieta's identities also give
\begin{equation}\label{eq:inverse-square}
 u:=\sum_{i=1}^3 f_i^{-2}=1-\frac{2b}{c}.
\end{equation}

\subsection{The signed quotient}

Let $\cS$ be the sphericalization of a fusion category $\cC$ with self-dual rank-three ring. There is a tensor functor $\pi:\cS\to\cC$. Each simple object $V$ of $\cC$ has two simple lifts $V_+,V_-$, both self-dual, with opposite dimensions. Put $\eps=\one_-$. Then
\[
 \eps^2\cong\one,\qquad \dim(\eps)=-1,\qquad
 \eps\otimes V_+\cong V_-.
\]
The construction and these properties are given in \cite[Remark~3.1]{ENO} and \cite[Section~5.1]{Ostrik}. A useful realization is the equivariantization for the canonical $C_2$-action \cite[Section~5.3]{Bartlett}.

All simples of $\cS$ are self-dual, so $K_0(\cS)$ is commutative. Its two quotients at $\eps=1$ and $\eps=-1$ are respectively
\[
 R=K_0(\cC),\qquad \widetilde R=K_0(\cS)/(\eps+1).
\]
Choose lifts of $X,Y$ and denote their images in $\widetilde R$ by $\widetilde X,\widetilde Y$. Then $\widetilde R=K(\widetilde k,\widetilde l,\widetilde m,\widetilde n)$, with the same multiplication formulas and associativity relation as above, but with integer, possibly negative, coefficients. For each corresponding pair of coefficients $a,\widetilde a$,
\begin{equation}\label{eq:signed}
 |\widetilde a|\leq a,\qquad \widetilde a\equiv a\pmod 2.
\end{equation}
Indeed, $a$ is the sum and $\widetilde a$ the difference of the two lift multiplicities. Changing the chosen lift changes the sign of the corresponding basis element.

The coefficient-of-$1$ functional $\tau$ on $\widetilde R$ satisfies
\begin{equation}\label{eq:orthonormal}
 \tau(\widetilde X_i\widetilde X_j)=\delta_{ij},
 \qquad (\widetilde X_0,\widetilde X_1,\widetilde X_2)=(1,\widetilde X,\widetilde Y).
\end{equation}
For example, $X_+^2$ contains $\one$ once and does not contain $\eps$, since $X_+$ is self-dual and $\eps X_+\not\cong X_+$. Thus taking the signed quotient preserves this pairing.

\begin{lemma}[Frobenius trace identity]\label{lem:trace}
Let $L$ be a number field with $\Q$-basis $e_0=1,e_1,\ldots,e_{r-1}$ and a linear functional $\tau$ satisfying $\tau(e_i e_j)=\delta_{ij}$. Put $\omega=\sum_i e_i^2$. If $\omega\ne0$, then
\[
 \Tr_{L/\Q}(a/\omega)=\tau(a)\qquad(a\in L).
\]
\end{lemma}
\begin{proof}
The $i$th diagonal entry of multiplication by $a$ is $\tau(e_i a e_i)$. Hence
$\Tr_{L/\Q}(a)=\tau(a\omega)$. Replace $a$ by $a/\omega$.
\end{proof}

\begin{lemma}\label{lem:field}
Suppose that $R$ is cubic. The dimension field $L$ of $\cS$ is a totally real cyclic cubic field, and the dimension homomorphism identifies $\widetilde R\otimes\Q$ with $L$. Set
\[
 x=\dim(X_+),\qquad y=\dim(Y_+),\qquad
 \delta=1+x^2+y^2=\dim(\cC),
\]
so that $\dim(\cS)=2\delta$. Then
\begin{equation}\label{eq:trace-relations}
 \Tr_{L/\Q}(1/\delta)=1,\qquad
 \Tr_{L/\Q}(x/\delta)=\Tr_{L/\Q}(y/\delta)=0.
\end{equation}
The formal codegrees of $\widetilde R$ are the three conjugates $g_i$ of $\delta$, counted with multiplicity. Write
\[
 \prod_i(t-g_i)=t^3-\widetilde b t^2+\widetilde c t-\widetilde c.
\]
Then
\begin{equation}\label{eq:comparison}
 0<\widetilde b\leq b,\qquad
 \frac{\widetilde c}{c}\in\Z_{>0},\qquad
 0<\frac{\widetilde b}{\widetilde c}\leq\frac bc.
\end{equation}
Moreover,
\begin{equation}\label{eq:v}
 v:=\sum_i g_i^{-2}=1-\frac{2\widetilde b}{\widetilde c}.
\end{equation}
\end{lemma}
\begin{proof}
The dimension homomorphism factors through $\widetilde R\otimes\Q$, so its image is a field of degree at most three. The formal codegrees of $\cS$ are
\begin{equation}\label{eq:six-codegrees}
 2f_1,2f_2,2f_3,2g_1,2g_2,2g_3:
\end{equation}
characters with $\eps=1$ come from $R$, and those with $\eps=-1$ come from $\widetilde R$; the two lifts double the sum defining each codegree. Every formal codegree of a spherical category lies in its dimension field \cite[Corollary~2.15]{Ostrik}. Therefore $L$ contains all three roots of the irreducible cubic $p$. It has degree exactly three and is its splitting field. Spherical dimensions are totally real \cite[Corollary~2.10]{ENO}, so $L$ is totally real. The map $\widetilde R\otimes\Q\to L$ is now an isomorphism.

Apply Lemma~\ref{lem:trace} to the basis $1,x,y$ and the functional \eqref{eq:orthonormal}. This gives \eqref{eq:trace-relations}. The embeddings of $L$ are precisely the three characters of the signed algebra, and its Casimir element has image $\delta$; hence its codegrees are $g_i=\sigma_i(\delta)>0$. Their reciprocal sum is one, giving the displayed polynomial and \eqref{eq:v}.

Formula \eqref{eq:b}, valid also for the signed coefficients, gives
\[
 \widetilde b=7+3\widetilde k^2+3\widetilde l^2+
                    \widetilde m^2+\widetilde n^2\leq b.
\]
Finally, formal codegrees divide the global dimension even without a pivotal structure \cite[Corollary~2.14]{Ostrik}. Thus $\delta/f_1$ is an algebraic integer in $L$. Taking its field norm yields
\[
 N_{L/\Q}(\delta/f_1)=\frac{\widetilde c}{c}\in\Z_{>0}.
\]
The last inequality in \eqref{eq:comparison} follows.
\end{proof}

\section{The rational and quadratic cases}\label{sec:quadratic}

The signed quotient also gives a proof of sphericality when the codegree
polynomial is reducible. We give the details because the main result should
not depend on a rank-three sphericality theorem. We use two general facts
about formal codegrees \cite[Proposition~2.2 and Corollary~2.14]{Ostrik}:
they divide the global dimension as algebraic integers, and they are
\emph{$d$-numbers}. In particular, if two Galois-conjugate codegrees are
the roots of $t^2-At+B$, then $A^2/B$ is an integer. This conclusion is
unchanged when both codegrees are multiplied by the same rational scalar.

For a signed quotient as in \eqref{eq:signed}, let
$\widetilde p(t)=t^3-\widetilde b t^2+\widetilde c t-\widetilde c$
be the characteristic polynomial of multiplication by
$1+\widetilde X^2+\widetilde Y^2$. The proof of the reciprocal-codegree
identity uses only the orthonormal pairing \eqref{eq:orthonormal}, and
therefore applies to this signed algebra. Its roots are positive real
numbers. Formula \eqref{eq:b} gives $\widetilde b\leq b$. Equality implies
that every signed coefficient has the same absolute value as its unsigned
counterpart. Signed associativity then forces
$\widetilde k\widetilde n=kn$ and $\widetilde l\widetilde m=lm$;
changing signs of $\widetilde X,\widetilde Y$ makes all four parameters
nonnegative. The corresponding simple lifts form a tensor subcategory
equivalent to $\cC$, so
\begin{equation}\label{eq:equal-b}
 \widetilde b=b\quad\Longrightarrow\quad \cC\text{ is spherical}.
\end{equation}

If $p$ splits over $\Q$, the Frobenius--Perron global dimension is one
of its integer roots. Hence $\cC$ is weakly integral and has its
canonical spherical structure by \cite[Propositions~8.23--8.24]{ENO}.
We turn to the remaining reducible case.

\begin{proposition}\label{prop:quadratic}
If $p$ has an irreducible quadratic factor, every categorification of
$K(k,l,m,n)$ admits a spherical structure.
\end{proposition}
\begin{proof}
We may assume the Frobenius--Perron global dimension is irrational,
as the integral case was just settled. Write
\begin{equation}\label{eq:qfactor}
 p(t)=(t-\gamma)(t^2-\alpha t+\beta),\qquad
 \gamma,\alpha,\beta\in\Z_{>0},\quad
 \beta=\frac{\alpha\gamma}{\gamma-1}.
\end{equation}
The quadratic roots are positive. The rational codegree $\gamma$ is at
least two; the Frobenius--Perron root is the larger quadratic root and
is strictly larger than $\gamma$. Hence $\alpha>\gamma$ and
$\beta>\gamma$.

Let $L$ be the dimension field of the sphericalization. Its degree is
at most three, and it contains both roots of the quadratic factor by
the formal-codegree field theorem used in Lemma~\ref{lem:field}.
Consequently $L$ is that quadratic splitting field. The dimension
character of $\widetilde R$ and its conjugate are distinct; its third
character is rational. Thus either all three roots of $\widetilde p$
are integers or
\begin{equation}\label{eq:signed-factor}
 \widetilde p(t)=(t-\widetilde\gamma)
       (t^2-\widetilde\alpha t+\widetilde\beta),
 \qquad
 \widetilde\beta=\frac{\widetilde\alpha\widetilde\gamma}
 {\widetilde\gamma-1},
\end{equation}
where $\dim(\cC)$ and its conjugate are the two roots of the quadratic
factor. The signed formal codegrees of $\cS$ are twice the roots of
$\widetilde p$; the original formal codegrees of $\cS$ are twice those
of $p$. General codegree divisibility and field norms therefore give
\begin{equation}\label{eq:quadratic-norm}
 \frac{\widetilde\beta}{\beta}
 =N_{L/\Q}\bigl(\dim(\cC)/f\bigr)\in\Z_{>0},
\end{equation}
where $f$ is either quadratic root of $p$ (and in the all-integral
signed case $\widetilde\beta$ denotes the square of
$\dim(\cC)$). We derive a contradiction from
$\widetilde b<b$.

First suppose all three signed codegrees are integers. The two
conjugate dimension characters have the same signed codegree, say
$\lambda$; call the third $\mu$. Both are positive integers and
$2/\lambda+1/\mu=1$. This leaves only
$(\lambda,\lambda,\mu)=(3,3,3)$ or $(4,4,2)$.
The first would give $\widetilde b=9$, impossible by
\eqref{eq:b} and signed associativity. In the second case
$\widetilde b=10$ and $\dim(\cC)=4$.
Codegree divisibility implies $\gamma\mid4$. The character
corresponding to the unique rational codegree $\gamma$ is fixed by
Galois, so its values on $X,Y$ are rational algebraic integers and
therefore integers $s,t$
on $X,Y$ with $\gamma=1+s^2+t^2$. Thus $\gamma=2$
(the case $\gamma=4$ would require $s^2+t^2=3$).
One of $s,t$ is zero. After exchanging $X,Y$, the fusion relations
give $K(k,l,m,n)=K(1,0,m,0)$.
The signed quotient has the same form with a coefficient
$\widetilde m$ satisfying $|\widetilde m|\leq m$.
The quadratic codegree product for $K(1,0,m,0)$ is
$2(m^2+8)$, and the analogous signed product is
$2(\widetilde m^2+8)$. The norm integrality in
\eqref{eq:quadratic-norm} forces
$(\widetilde m^2+8)/(m^2+8)=1$.
Choose the sign of the lift so that $\widetilde m=m$.
Its products close in the sphericalization, proving sphericality.

Now assume \eqref{eq:signed-factor}. Put
$h=\widetilde\beta/\beta\in\Z_{>0}$.
If $\widetilde b=b$, we are done by \eqref{eq:equal-b}; otherwise
$\widetilde b<b$. Since $\widetilde\gamma\geq2$ and
$\widetilde\alpha+\widetilde\gamma<\alpha+\gamma$,
\[
 h\beta=\widetilde\beta
 =\widetilde\alpha\frac{\widetilde\gamma}
                      {\widetilde\gamma-1}
 <2(\alpha+\gamma-\widetilde\gamma)<4\beta.
\]
Here the last inequality follows from $\beta>\gamma$.
So $h\in\{1,2,3\}$. The rational codegree $\gamma$ divides
$\dim(\cC)$; taking norms shows
$h\beta/\gamma^2=h\alpha/[\gamma(\gamma-1)]$
is a positive integer. If $h=2$ or $3$, write this integer as $q$.
As $\widetilde\beta<2\widetilde\alpha$,
\[
 \widetilde\alpha\geq\frac h2\beta
 =\frac h2\alpha+\frac q2\gamma.
\]
For $h=3$ and $q\geq2$ the right side exceeds
$\alpha+\gamma$; for $h=2$ and $q\geq2$ it is at least
$\alpha+\gamma$. Both contradict $\widetilde b<b$.
If $q=1$, then
$\alpha=\gamma(\gamma-1)/h$ and
$\beta=\gamma^2/h$. As $\beta$ is integral, $h=2$ forces
$\gamma$ even and $h=3$ forces $\gamma$ divisible by three.
Consequently $\alpha^2/\beta=(\gamma-1)^2/h$
is not an integer, contrary to the $d$-number property of the
quadratic original codegrees. Hence $h=1$.

When $\widetilde\beta=\beta$, the inequality
$\widetilde b<b$ forces $\widetilde\gamma<\gamma$:
otherwise $\widetilde\alpha\geq\alpha$ by
\eqref{eq:qfactor}--\eqref{eq:signed-factor}.
The original and signed quadratic factors have the same splitting
field $L$. By the $d$-number property, both normalized
discriminants below are positive integers. The discriminants
$\alpha^2-4\beta$ and $\widetilde\alpha^2-4\beta$ generate
the same quadratic field, so their ratio is a rational square.
Dividing each by the same $\beta$ preserves that ratio; hence
the normalized integer radicands have the same squarefree part.
There are therefore a positive squarefree integer $d$
and positive integers $x,y$ with
\[
 \frac{\alpha^2}{\beta}-4=dx^2,\qquad
 \frac{\widetilde\alpha^2}{\beta}-4=dy^2.
\]
The ratio of the left sides after adding $4$ is the square of
\[
 \frac{\alpha}{\widetilde\alpha}
 =\frac{\widetilde\gamma/(\widetilde\gamma-1)}
        {\gamma/(\gamma-1)}=\frac AB,
 \qquad 1<\frac AB<2,
\]
where $A,B$ are coprime positive integers. Thus
$d(Bx-Ay)(Bx+Ay)=4(A^2-B^2)$.
The two factors on the left have the same parity and must both be
even. Put $q=(Bx-Ay)/2>0$. Then
\begin{equation}\label{eq:quad-descent}
 Bx=\frac{A^2-B^2+dq^2}{dq},\qquad
 Ay=\frac{A^2-B^2-dq^2}{dq}.
\end{equation}
Reducing the first equality modulo $B$ and the second modulo $A$
shows that $A^2+B^2+dq^2$ is divisible by $AB$.
Since $Ay>0$, we have $dq^2<A^2-B^2$, and hence
\[
 2<\frac{A^2+B^2+dq^2}{AB}<\frac{2A}{B}<4.
\]
The integer in the middle is therefore three. Substitution into
\eqref{eq:quad-descent} gives
\[
 x=\frac{3A-2B}{dq},\qquad
 y=\frac{2A-3B}{dq}.
\]
Because $\gcd(A,B)=1$, these equations imply $dq\mid5$.
Thus $dq^2$ is $1$, $5$ or $25$.
If $dq^2=1$, the equation
$A^2+B^2+1=3AB$ divided by $B^2$ gives
$B^{-2}=3(A/B)-(A/B)^2-1>1$, impossible.
If $dq^2=5$, the same inequality gives $B^2<5$, so
$B=1$ or $2$; the former allows no integer $A/B\in(1,2)$,
whereas the latter forces $(A,B)=(3,2)$ and then $y=0$.
Finally, if $dq^2=25$, completing the square gives
$(2A-3B)^2-5B^2=-100$; reduction modulo five and then
modulo twenty-five forces both $A,B$ divisible by five.
All cases are impossible. We conclude $\widetilde b=b$ and
\eqref{eq:equal-b} gives sphericality.
\end{proof}

\section{Central twists and the sign object}

For a spherical fusion category $\cD$, let $I:\cD\to Z(\cD)$ be the right adjoint to the forgetful functor. Write $d=\dim(\cD)$, and for $A\in\Irr Z(\cD)$ put $a_A=\dim(A)$ and let $\theta_A$ be its twist. The numbers $a_A$ are real, and the twists are roots of unity. We use the identities
\begin{align}
 \sum_A a_A^2&=d^2,\label{eq:mass}\\
 \sum_A a_A^2\theta_A^n
   &=d\sum_{V\in\Irr\cD}\dim(V)\nu_n(V)\qquad(n\geq1),\label{eq:moments}
\end{align}
where $\nu_n$ is the $n$th Frobenius--Schur indicator for the specified pivotal structure. To see \eqref{eq:moments}, multiply the Ng--Schauenburg identity
\[
 \tr(\theta^n|_{I(V)})=d\nu_n(V)
\]
by $\dim(V)$ and sum over $V$. Adjunction and sphericality of the forgetful functor give
$a_A=\sum_V[I(V):A]\dim(V)$, yielding the formula. The indicator identity is \cite[Theorem~4.1]{NSexp}; \eqref{eq:mass} is the global-dimension formula for the center. Also,
\begin{equation}\label{eq:cyclo}
 \nu_n(V)\in\Q(\zeta_n),
\end{equation}
because the cyclic operator defining the indicator has $n$th power equal to the identity \cite[Theorem~5.1]{NSind}. These facts require sphericality, but not pseudo-unitarity.

\begin{lemma}[The central sign]\label{lem:sign}
The sign object of a sphericalization has a canonical lift, again denoted $\eps$, to $Z(\cS)$, satisfying
\[
 \eps^2\cong\one,\qquad c_{\eps,\eps}=\id,
 \qquad \dim(\eps)=-1,\qquad \theta_\eps=-1.
\]
For its underlying object in $\cS$,
\begin{equation}\label{eq:eps-ind}
 \nu_n(\eps)=
 \begin{cases}0,&n\text{ odd},\\-1,&n\text{ even}.
 \end{cases}
\end{equation}
\end{lemma}
\begin{proof}
View $\cS$ as the equivariantization for the canonical $C_2$-action \cite[Definition~5.13 and Corollary~5.15]{Bartlett}. The objects supported on the tensor unit form the canonical central copy of $\Rep(C_2)$. The sign representation has the usual symmetric self-braiding, and its dimension in the canonical spherical structure is $-1$. On this copy of $\Rep(C_2)$ the Drinfeld isomorphism is the usual one, whereas the pivotal map of the sign representation is minus the usual map; hence its twist is $-1$.

The space $\Hom(\one,\eps^{\otimes n})$ is zero for odd $n$ and one-dimensional for even $n$. In the latter case the ordinary cyclic rotation for $\Rep(C_2)$ is the identity. The pivotal structure used here introduces one inverse pivotal map of $\eps$, namely the scalar $-1$, into that rotation. Its trace is therefore $-1$, proving \eqref{eq:eps-ind}.
\end{proof}

\begin{lemma}[A forced squared-twist sector]\label{lem:sector}
In the cubic situation of Lemma~\ref{lem:field}, the simple constituents of $I(\one)$ and $I(\eps)$ are twelve distinct simple objects of $Z(\cS)$. Their squared twists are all one, and the sum of their squared dimensions is
\begin{equation}\label{eq:known-mass}
 2\delta^2(u+v).
\end{equation}
The same statement holds after every embedding of $L$.
\end{lemma}
\begin{proof}
Since $K_0(\cS)$ is commutative, $I(\one)$ is multiplicity-free. The formal-codegree theorem \cite[Theorem~2.13]{Ostrik} and \eqref{eq:six-codegrees} give its six simple constituent dimensions as
\[
 \frac{\delta}{f_1},\frac{\delta}{f_2},\frac{\delta}{f_3},
 \frac{\delta}{g_1},\frac{\delta}{g_2},\frac{\delta}{g_3}.
\]
They are positive under every real embedding and their squared sum is $\delta^2(u+v)$. Their twists are one \cite[Theorem~2.5]{Ostrik}.

The projection formula gives $I(\eps)\cong\eps\otimes I(\one)$. Thus its six constituents have the opposite, negative dimensions, and are disjoint from those of $I(\one)$. To check their squared twists, let $\lambda_A$ be the scalar monodromy of $\eps$ with a simple object $A$. The normalized $S$-matrix column $B\mapsto s_{A,B}/a_A$ is a ring character of $K_0(Z(\cS))$ \cite[Example~2.9]{Ostrik}. Its value on $\eps$ is $\dim(\eps)\lambda_A=-\lambda_A$. Since $\eps^2=\one$, it follows that $\lambda_A^2=1$. Balancing and $\theta_\eps=-1$ now give $\theta_{\eps\otimes A}^2=1$ whenever $\theta_A=1$. The dimensions are negated, so their squares give a second copy of the same sum. This proves \eqref{eq:known-mass}.
\end{proof}

\section{Galois averaging and the cubic bound}

\begin{proposition}[Cubic codegree bound]\label{prop:bound}
If $K(k,l,m,n)$ is cubic and admits a complex fusion categorification, then
\begin{equation}\label{eq:bound}
 c\leq\frac72 b.
\end{equation}
\end{proposition}
\begin{proof}
Use the notation of Lemmas~\ref{lem:field}--\ref{lem:sector}, and put
\[
 M=\Q(\zeta_{120}).
\]
The fields $L$ and $M$ are linearly disjoint: $L/\Q$ is Galois of degree $3$, whereas $[M:\Q]=\varphi(120)=32$. Choose three automorphisms $\sigma_1,\sigma_2,\sigma_3$ of $\overline\Q$ restricting to the three embeddings of $L$ and fixing $M$ pointwise. Fix these choices for the rest of the proof. In particular, they fix every indicator of each of the orders
\[
 2,4,6,8,10,12,
\]
by \eqref{eq:cyclo}. The same choices work for all six orders, since each divides $120$. We apply these automorphisms only to algebraic scalar identities; no choice of a field of definition for the associators is needed.

On the unit circle $\T$, define the finite positive measure
\begin{equation}\label{eq:measure}
 \mu=\sum_{i=1}^3\ \sum_{A\in\Irr Z(\cS)}
       \frac{\sigma_i(a_A)^2}{4\sigma_i(\delta)^2}
       \Dirac_{\sigma_i(\theta_A)^2}.
\end{equation}
Here $\Dirac_z$ is unit point mass at $z$. Every central dimension belongs to $L$, since the forgetful functor is spherical. Its conjugates are real, and $\sigma_i(\delta)>0$. Thus the coefficients in \eqref{eq:measure} are nonnegative. Twists and their conjugates are roots of unity, so the support is contained in $\T$. Equation~\eqref{eq:mass} gives
\begin{equation}\label{eq:m0}
 \mu(\T)=3.
\end{equation}

For $1\leq j\leq6$, apply $\sigma_i$ to \eqref{eq:moments}, with $n=2j$, and divide by $4\sigma_i(\delta)^2$. Since the indicators are fixed, summing over $i$ gives
\begin{align}
 \int_{\T}z^j\,d\mu(z)
 &=\frac12\sum_{V\in\Irr\cS}\nu_{2j}(V)
       \Tr_{L/\Q}\left(\frac{\dim(V)}{\delta}\right)\notag\\
 &=\frac12\bigl(\nu_{2j}(\one)-\nu_{2j}(\eps)\bigr)=1.
 \label{eq:averaged}
\end{align}
For the second equality, the four non-unit lifts have dimensions $\pm x,\pm y$ and vanish under the weighted trace by \eqref{eq:trace-relations}; the traces for $\one$ and $\eps$ are $1$ and $-1$. The last equality uses Lemma~\ref{lem:sign} and $\nu_n(\one)=1$. This step does not require knowing any indicator of $X_\pm$ or $Y_\pm$.

Lemma~\ref{lem:sector} forces the lower bound
\begin{equation}\label{eq:atom-lower}
 \mu(\{1\})\geq\frac32(u+v).
\end{equation}
Indeed, in each of the three summands defining $\mu$, those twelve objects contribute $(u+v)/2$ at one. The sums $u,v$ are rational and hence independent of the embedding.

Now use the nonnegative trigonometric polynomial
\[
 P(z)=|1+z+\cdots+z^6|^2
     =7+\sum_{j=1}^6(7-j)(z^j+z^{-j})\qquad(z\in\T).
\]
Its value at one is $49$. Since the negative moments are the complex conjugates of the positive ones, \eqref{eq:m0} and \eqref{eq:averaged} imply
\[
 49\mu(\{1\})\leq\int_{\T}P(z)\,d\mu(z)
   =7\cdot3+2\sum_{j=1}^6(7-j)=63.
\]
Combining this with \eqref{eq:atom-lower} yields
\[
 u+v\leq\frac67.
\]
Substitute \eqref{eq:inverse-square} and \eqref{eq:v} to obtain
\begin{equation}\label{eq:two-ratios}
 \frac bc+\frac{\widetilde b}{\widetilde c}\geq\frac47.
\end{equation}
But \eqref{eq:comparison} bounds the second ratio by the first. Therefore $b/c\geq2/7$, which is \eqref{eq:bound}.
\end{proof}

\begin{remark}[The role of the six indicators]\label{rem:six}
The moment argument uses only the fields containing the indicators, not their values, signs or sizes. Its atom bound is sharp for the stated moment information: the measure with mass $9/7$ at $1$ and mass $2/7$ at each nontrivial seventh root has mass $3$ and first six moments $1$. The choice of order twelve is also natural for field disjointness. At order fourteen, $\Q(\zeta_{14})$ already contains a cubic field, so the degree argument used above cannot be continued uniformly.
\end{remark}

\section{Reducible rings: an independent exclusion}
\label{sec:reducible-rings}

We complete the classification of \emph{fusion rings} without invoking
the rank-three classification theorem of \cite{Ostrik}. The cubic rings
were handled in Proposition~\ref{prop:cubic-classification}; in this
section $p$ is reducible. We use the general formal-codegree theorem
for the simple summands of $I(\one)$ and the two center identities
\eqref{eq:mass} and \eqref{eq:moments}.

\begin{lemma}[Gauss bound]\label{lem:gauss-bound}
If a reducible rank-three ring $K(k,l,m,n)$ admits a complex fusion
categorification, and $f_{\max}$ is its largest codegree, then
\begin{equation}\label{eq:gauss-bound}
 1-\frac{4b}{c}\leq\frac{1}{f_{\max}}.
\end{equation}
Consequently, if $b\geq50$, a categorification requires $c-4b\leq4$.
\end{lemma}
\begin{proof}
The category is spherical by Theorem~\ref{thm:main}. Choose its
canonical positive spherical structure when its Frobenius--Perron
global dimension is an integer \cite[Propositions~8.23--8.24]{ENO}.
Otherwise $p$ has an irreducible quadratic factor. Since $\cC$ is
spherical, its formal codegrees belong to the field generated by
its simple dimensions \cite[Corollary~2.15]{Ostrik}. The latter
is generated by a ring character of a three-dimensional
$\Q$-algebra, so has degree at most three. Since it contains the
quadratic splitting field of $p$, it is exactly that quadratic field;
in particular the dimension character is one of the two quadratic
characters. The positive Frobenius--Perron character is one of
these two conjugates. Apply a field automorphism
to scalar identities so that the simple dimensions become their
positive Frobenius--Perron values. Every center simple dimension is
the dimension of its image under the forgetful functor, a nonnegative
integral combination of these values, and hence is positive. Twists
remain roots of unity. Write $D=f_{\max}$ for the resulting global
dimension.

Because the ring is commutative, $I(\one)$ has three distinct simple
summands with dimensions $D/f_1,D/f_2,D/f_3$ and trivial twists
\cite[Theorems~2.5 and 2.13]{Ostrik}. Let
$H=D^2\sum_i f_i^{-2}$ be their total squared dimension.
The remaining center simples have total squared dimension $D^2-H$
by \eqref{eq:mass}, and their twist-weighted sum is $D-H$ by
\eqref{eq:moments} with $n=1$, since $\nu_1(\one)=1$ and all
other first indicators vanish. The triangle inequality gives
$|D-H|\leq D^2-H$, hence
$\sum_i f_i^{-2}\leq(1+D^{-1})/2$.
Since $\sum_i f_i^{-2}=1-2b/c$ by \eqref{eq:inverse-square},
this is \eqref{eq:gauss-bound}.

Suppose $c-4b>4$ and $b\geq50$. The bound implies
$f_{\max}\leq c/(c-4b)$. Since $f_{\max}\geq b/3>9$,
we get $c<\tfrac92 b$.
Now $p(1)=1-b<0$, $p(2)=8+c-4b>0$ and $p(b)>0$.
Also
\[
 p(b-5)=(b-5)\left(c-\frac92b+\frac{50-b}{2}\right)-c<0.
\]
Thus $f_{\max}>b-5$. As $c-4b$ is an integer at least five,
\[
 b-5<f_{\max}\leq\frac{c}{c-4b}\leq\frac c5,
\]
so $c>5b-25\geq\tfrac92b$, a contradiction.
\end{proof}

\begin{lemma}[Arithmetic reduction]\label{lem:arithmetic-reduction}
If a reducible ring $K(k,l,m,n)$ admits a complex categorification,
then it is either $K(1,0,m,0)$ for some $m\geq0$, or $K(2,1,2,1)$.
\end{lemma}
\begin{proof}
If one of $k,l$ is zero, interchange $X,Y$ to arrange $l=0$.
Then
\eqref{eq:assoc} gives $k(k-n)=1$, so $k=1,n=0$ and we obtain
$K(1,0,m,0)$.

Assume $k,l>0$. When $b<50$, equation \eqref{eq:b} gives
$k,l\leq3$ and $m,n\leq6$. Checking these four finite ranges against
\eqref{eq:assoc} and
$7+3(k^2+l^2)+m^2+n^2<50$ leaves the following tuples,
up to exchanging $X,Y$:
\[
 \begin{array}{c|c|c}
 (k,l,m,n)&(b,c)&\text{factorization or irreducibility check}\\\hline
 (1,1,1,0)&(14,49)&\text{irreducible modulo }2\\
 (2,1,0,2)&(26,148)&\text{irreducible modulo }3\\
 (2,1,2,1)&(27,108)&(t-3)(t^2-24t+36)\\
 (2,1,4,0)&(38,148)&\text{irreducible modulo }3\\
 (3,1,0,3)&(46,473)&\text{irreducible modulo }2
 \end{array}
\]
For the two rows with $c=148$, reduction of
$t^3-bt^2+ct-c$ modulo three is $t^3+t^2+t+2$,
which has no root in $\mathbb F_3$.
The first and last rows have reduction $t^3+t+1$ modulo two,
with no root in $\mathbb F_2$. The first row is excluded from
consideration here because $p$ is
assumed reducible. Thus the only candidate with $l>0$ and $b<50$
is $K(2,1,2,1)$.

It remains to consider $b\geq50$. Lemma~\ref{lem:gauss-bound}
gives $c-4b\leq4$. We retain the reducibility assumption:
the small value of $c-4b$ alone also admits cubic Pell-type
families, handled in Proposition~\ref{prop:cubic-classification}. Set
\[
 A=kl-mn,\qquad B=2(k^2-lm)-1.
\]
By \eqref{eq:c}, $A^2+B^2\leq13$. Exchanging $X,Y$ changes
$B$ to $-B$ (use \eqref{eq:assoc}), so choose $B>0$.
As $B$ is odd, it equals $1$ or $3$.

If $B=1$, then $k^2-lm=1$ and $l^2-kn=0$.
Since $\gcd(k,l)=1$ by \eqref{eq:assoc}, $k\mid l^2$
forces $k=1$; hence $n=l^2$ and $lm=0$.
For $l=0$ this is the near-group family. For $l>0$ we have
$m=0$ and the character value of $X$ satisfies
\[
 q_l(z)=z^3-lz^2-2z+l=0.
\]
This monic cubic has no integer root for $l\geq1$:
$q_l(0),q_l(1),q_l(-1)$ are nonzero, and any other integer root $r$
would imply $(r^2-1)\mid r(r^2-2)$, hence $(r^2-1)\mid r$,
which is impossible. Hence $\Q[X]$ is a cubic number field.
The element $\Omega$ is nonrational, since its trace is $b>9$
whereas a rational $\Omega$ would have three equal codegrees, all
equal to three. A cubic field has no proper intermediate fields,
so $\Omega$ also generates $\Q[X]$; in particular $p$ is
irreducible, contrary to our reducibility assumption.

If $B=3$, the relation $k^2-lm=2$ gives $l^2-kn=-1$;
in particular $k,l>0$ and
\[
 m=\frac{k^2-2}{l},\quad n=\frac{l^2+1}{k},\quad
 A=\frac{2l^2-k^2+2}{kl}\in\{-2,-1,0,1,2\}.
\]
For $A=\pm2$ one gets $3l^2+2=(k\pm l)^2$,
impossible modulo three. For $A=\pm1$ one gets
$9l^2+8=(2k\pm l)^2$, equally impossible modulo three.
Thus $A=0$, so
\[
 k^2=2l^2+2,\qquad m=2l,\qquad n=k/2,\qquad
 b=\frac{27}{2}(l^2+1),\qquad c=4b.
\]
Here $k$ is even, $l$ is odd and $b$ is odd. If $l=1$,
this is $K(2,1,2,1)$. Otherwise $b\geq135$ and
$p(t)=t^3-b(t-2)^2$ is irreducible: an integer root $r$
would satisfy $r-2\mid p(2)=8$ and
$b=r^3/(r-2)^2$; the possible positive roots
$r=3,4,6,10$ give respectively $b=27,16,27/2,125/8$.
This completes the arithmetic reduction.
\end{proof}

\begin{proposition}[A uniform near-group exclusion]
\label{prop:near-group}
If $K(1,0,m,0)$ admits a complex fusion categorification, then
$m\leq2$.
\end{proposition}
\begin{proof}
Write $Y^2=\one$, $XY=X$ and $X^2=\one\oplus Y\oplus mX$.
Assume $m>2$ and set $d=\FPdim(X)=(m+\sqrt{m^2+8})/2$.
The category is spherical by Theorem~\ref{thm:main}. Its dimension
character is one of the two quadratic ring characters: indeed
the irrational formal codegrees belong to its dimension field by
the general field theorem. Apply an automorphism to all algebraic
scalar identities so that the simple dimensions become
$\dim(\one)=\dim(Y)=1$, $\dim(X)=d>0$.
Write $\Delta=2+d^2$; every central simple dimension is positive
under this embedding, and all conjugated twists have modulus one.

The rank-three formal-codegree decomposition of $I(\one)$ gives
$I(\one)=\one\oplus A\oplus B$, with $\dim(B)=\Delta/2$ and
$\dim(A)=\Delta/2-1$. Since
$FI(\one)=3\one\oplus Y\oplus mX$ and both $F(A),F(B)$
contain $\one$ once, comparison of rational and $d$ coefficients
gives $m=2t$ for an integer $t\geq2$ and
\begin{equation}\label{eq:near-I1}
 F(A)=\one\oplus tX,\quad F(B)=\one\oplus Y\oplus tX,
 \quad\theta_A=\theta_B=1.
\end{equation}
The irrationality of $d$ for $m>2$ justifies the coefficient
comparison. Hence an odd $m>2$ is already impossible.

Adjunction, applied to
$FI(Y)=\one\oplus3Y\oplus2tX$, gives
\begin{equation}\label{eq:near-IY}
 I(Y)=B\oplus C_1\oplus C_2,\qquad
 F(C_1)=Y\oplus\alpha X,\quad
 F(C_2)=Y\oplus\beta X,\quad \alpha+\beta=t.
\end{equation}
Indeed, $\dim\operatorname{End}I(Y)=3$, and $B$ is its sole
summand from $I(\one)$, with multiplicity one.
The first twist trace $\operatorname{Tr}(\theta|_{I(Y)})=0$
follows from the Ng--Schauenburg identity used in \eqref{eq:moments},
Section~4.
As $\dim(B)=\dim(C_1)+\dim(C_2)>0$ and $\theta_B=1$,
equality in the triangle inequality forces
$\theta_{C_1}=\theta_{C_2}=-1$.

The remaining central simples in $I(X)$ have the form $E_j$
with $F(E_j)=r_jX$ and multiplicity $r_j>0$.
By adjunction, using \eqref{eq:near-I1}--\eqref{eq:near-IY},
\[
 I(X)=tA\oplus tB\oplus\alpha C_1\oplus\beta C_2
             \oplus\bigoplus_j r_jE_j.
\]
Since $FI(X)=2t(\one\oplus Y)\oplus(4t^2+4)X$,
the dimension of $\operatorname{End}I(X)$ is $4t^2+4$.
Taking the sum of squared multiplicities and writing
$R=\sqrt{t^2+2}$ gives
\begin{equation}\label{eq:near-mass}
 M:=\sum_j r_j^2=t^2+4+2\alpha\beta.
\end{equation}
The Ng--Schauenburg twist trace for $I(X)$ and its second-power
version give, with $z_j=\theta_{E_j}$ and the ordinary
second indicator $\nu_2(X)=\eta\in\{1,-1\}$,
\begin{align}
 S_1:=\sum_j r_j^2z_j
   &=-2\alpha\beta-tR,\label{eq:near-first}\\
 S_2:=\sum_j r_j^2z_j^2
   &=2\alpha\beta-t^2-(2t-2\eta)R.
   \label{eq:near-second}
\end{align}
These equations follow by expanding the known summands of $I(X)$
in the twist-trace identity for $q=1,2$:
their dimensions are $1+td$, $2+td$, $1+\alpha d$,
$1+\beta d$, with twists $1,1,-1,-1$, and
$d=t+R$, $d^{-1}=(R-t)/2$.

The polynomial $P(z)=|1+z+z^2|^2$ is nonnegative for every
$|z|=1$. Thus \eqref{eq:near-mass}--\eqref{eq:near-second}
imply
\[
 0\leq\sum_j r_j^2P(z_j)
   =3M+4S_1+2S_2
   =t^2+12+2\alpha\beta-(8t-4\eta)R.
\]
But $\alpha\beta\leq t^2/4$, $\eta\leq1$, and $R>t$.
The last expression is strictly less than
$\tfrac32t^2+12-(8t-4)t
=12+4t-\tfrac{13}{2}t^2<0$ for every $t\geq2$.
This contradiction excludes all $m>2$ at once.
\end{proof}

\begin{proposition}[The quadratic exceptional ring]
\label{prop:quadratic-exception}
The ring $K(2,1,2,1)$ has no complex fusion categorification.
\end{proposition}
\begin{proof}
Suppose it has a categorification. By Theorem~\ref{thm:main}
it is spherical, and the field-of-dimensions argument in
Proposition~\ref{prop:quadratic} lets us apply a field automorphism
to scalar identities so that
\[
 x:=\dim(X)=2+\sqrt3,\qquad y:=\dim(Y)=1+\sqrt3,
 \qquad \Delta=12+6\sqrt3.
\]
All central simple dimensions then become positive. The formal
codegrees are $\Delta,3,12-6\sqrt3$, so
$I(\one)=\one\oplus A\oplus B$ with
$\dim(A)=\Delta/3=4+2\sqrt3$ and
$\dim(B)=\Delta/(12-6\sqrt3)=7+4\sqrt3$.
Comparing $FI(\one)=3\one\oplus3X\oplus3Y$ with these
dimensions and using $[F(A):\one]=[F(B):\one]=1$ gives
\[
 F(A)=\one\oplus X\oplus Y,
 \qquad F(B)=\one\oplus2X\oplus2Y,
 \qquad \theta_A=\theta_B=1.
\]
Adjunction now gives
$I(Y)=A\oplus2B\oplus\bigoplus_j s_jS_j$, where the
$S_j$ are center simples distinct from $\one,A,B$,
$FI(Y)=3\one\oplus9X\oplus9Y$ and
$\sum_j s_jF(S_j)=4X\oplus4Y$. The twist-trace identity says
$\operatorname{Tr}(\theta|_{I(Y)})=0$. Because twists have modulus
one and all conjugated dimensions are positive, this would force
\[
 \dim(A)+2\dim(B)\leq\sum_j s_j\dim(S_j),
 \qquad 18+10\sqrt3\leq12+8\sqrt3,
\]
which is false.
\end{proof}

\section{Completion of the classification}

\begin{proposition}\label{prop:cubic-classification}
A cubic rank-three fusion ring admits a complex fusion categorification only if it is isomorphic to $K(1,1,1,0)$. Every categorification of this ring admits a spherical structure.
\end{proposition}
\begin{proof}
Combine \eqref{eq:lower} and Proposition~\ref{prop:bound}:
\[
 4b-8\leq c\leq\frac72b,
 \qquad\text{hence}\qquad b\leq16.
\]
If one of $k,l$ is zero, interchange $X,Y$ so that $l=0$. Equation~\eqref{eq:assoc} becomes $k(k-n)=1$, giving $k=1,n=0$. The character $X\mapsto0$, $Y\mapsto-1$ then has formal codegree $2$, contrary to irreducibility of $p$. Thus $k,l\geq1$.

By \eqref{eq:b},
\[
 3(k^2+l^2)+m^2+n^2\leq9.
\]
This forces $k=l=1$. Associativity now gives $m+n=1$. The two resulting choices are interchanged by exchanging $X$ and $Y$, so the ring is $K(1,1,1,0)$.

It remains to descend sphericality to the given categorification, rather than merely exhibit a different spherical categorification. For $K(1,1,1,0)$, \eqref{eq:signed} gives
\[
 \widetilde k,\widetilde l,\widetilde m\in\{1,-1\},
 \qquad\widetilde n=0.
\]
Signed associativity implies $\widetilde l\widetilde m=1$. Replacing $\widetilde X$ by $\widetilde l\widetilde X$ and $\widetilde Y$ by $\widetilde k\widetilde Y$ makes all four signed parameters equal to the original ones. For these choices of lifts, each signed multiplicity equals the corresponding unsigned multiplicity. Since these are respectively the difference and sum of the two lift multiplicities, the unwanted lift has multiplicity zero in every product. The three chosen lifts therefore form a fusion subcategory of $\cS$.

The restriction of $\pi$ to this subcategory is a tensor equivalence with $\cC$: it is bijective on simple objects and fully faithful. Transporting its spherical structure gives a spherical structure on $\cC$. This is also the section criterion of \cite[Section~5.1(c)]{Ostrik}.
\end{proof}

\begin{proof}[Proof of Theorem~\ref{thm:main}]
If the Grothendieck ring is $\Z[C_3]$, the category is pointed and hence admits a spherical structure. Otherwise it is $K(k,l,m,n)$. The rational and quadratic cases are proved in Section~\ref{sec:quadratic}; the irreducible cubic case is Proposition~\ref{prop:cubic-classification}. These cases exhaust the characteristic polynomial of degree three.
\end{proof}

\begin{corollary}\label{cor:list}
Up to based-ring isomorphism, the rank-three fusion rings admitting a fusion categorification over $\C$ are precisely
\[
 \Z[C_3],\quad K(1,0,0,0),\quad K(1,0,1,0),\quad
 K(1,0,2,0),\quad K(1,1,1,0).
\]
Each has a unitary realization. Thus general, pivotal, spherical, pseudo-unitary and unitary \emph{existence of a categorification} give the same list in rank three, with no multiplicity restriction.
\end{corollary}
\begin{proof}
For necessity, the pointed case is $\Z[C_3]$. A nonpointed ring
with irreducible cubic codegree polynomial is $K(1,1,1,0)$ by
Proposition~\ref{prop:cubic-classification}. By
Lemma~\ref{lem:arithmetic-reduction}, a reducible categorifiable ring
is $K(2,1,2,1)$ or $K(1,0,m,0)$. The first is excluded by
Proposition~\ref{prop:quadratic-exception}, and
Proposition~\ref{prop:near-group} forces $m\leq2$.
Conversely, the following standard unitary categories realize the five rings:
\begin{center}
\begin{tabular}{@{}ll@{}}
\toprule
Fusion ring & A unitary realization\\
\midrule
$\Z[C_3]$ & $\Vect_{C_3}$\\
$K(1,0,0,0)$ & An Ising category\\
$K(1,0,1,0)$ & $\Rep(S_3)$\\
$K(1,0,2,0)$ & The rank-three even part of the $E_6$ subfactor\\
$K(1,1,1,0)$ & The even subcategory of $SU(2)_5$\\
\bottomrule
\end{tabular}
\end{center}
These are the realizations in \cite{Ostrik}; see also \cite{HH} for the $E_6$ fusion rules. Unitary categories are pseudo-unitary and spherical, so the existence statements coincide.
\end{proof}

\begin{remark}
The corollary concerns five \emph{fusion rings}, not five
tensor-equivalence classes of categories. The full category list,
including the pointed twists, quantum-group choices and
Galois-conjugate realizations, remains the one in
\cite[Theorem~1.1]{Ostrik}, now without its pivotal hypothesis.
No assertion is made that every categorification on the list is unitary.
\end{remark}

\paragraph{AI disclosure.}
GPT-6 Astra assisted with mathematical exploration and manuscript preparation.
This manuscript and its accompanying research data are AI-generated. They should be treated with caution: they have not been independently verified by a human mathematician or formalized in Lean.

\begin{thebibliography}{99}
\bibitem{Bartlett}
B.~Bartlett, \emph{Fusion categories via string diagrams}, preprint (2015),
\href{https://arxiv.org/abs/1502.02882}{arXiv:1502.02882v2}.

\bibitem{ENO}
P.~Etingof, D.~Nikshych and V.~Ostrik, \emph{On fusion categories},
Ann. of Math. (2) \textbf{162} (2005), no.~2, 581--642;
\href{https://arxiv.org/abs/math/0203060}{arXiv:math/0203060} (updated version).

\bibitem{HH}
T.~J.~Hagge and S.-M.~Hong, \emph{Some non-braided fusion categories of rank three},
Commun. Contemp. Math. \textbf{11} (2009), no.~4, 615--637;
\href{https://arxiv.org/abs/0704.0208}{arXiv:0704.0208}.

\bibitem{NSind}
S.-H.~Ng and P.~Schauenburg, \emph{Higher Frobenius--Schur indicators for pivotal categories},
in \emph{Hopf Algebras and Generalizations}, Contemp. Math. \textbf{441},
Amer. Math. Soc., Providence, RI, 2007, 63--90;
\href{https://doi.org/10.1090/conm/441/08500}{doi:10.1090/conm/441/08500}.

\bibitem{NSexp}
S.-H.~Ng and P.~Schauenburg, \emph{Frobenius--Schur indicators and exponents of spherical categories},
Adv. Math. \textbf{211} (2007), no.~1, 34--71;
\href{https://doi.org/10.1016/j.aim.2006.07.017}{doi:10.1016/j.aim.2006.07.017}.

\bibitem{Ostrik}
V.~Ostrik, \emph{Pivotal fusion categories of rank 3}, with an appendix written jointly with D.~Nikshych,
Mosc. Math. J. \textbf{15} (2015), no.~2, 373--396;
\href{https://arxiv.org/abs/1309.4822}{arXiv:1309.4822}.
\end{thebibliography}
\end{document}

